\section{Descrizione Generale}

\subsection{Obiettivo del prodotto}
Il prodotto prevede lo sviluppo incrementale (attraverso le fasi di Proof of Concept\textsubscript{G} e Minimum Viable Product\textsubscript{G}) di un'applicazione web per la redazione di testi in formato Markdown. 

L'obiettivo principale è superare i limiti dei classici chatbot\textsubscript{G} testuali, fornendo un ambiente editor integrato in cui l'LLM possa attivamente elaborare, criticare e generare parti di testo su richiesta dell'utente, migliorando la produttività. 

L'obiettivo a lungo termine (vincolato ai requisiti opzionali) prevede la creazione di una knowledge base strutturata in un database server-side, che permetta il collegamento logico delle note per fornire un contesto più accurato e informato alle risposte dell'intelligenza artificiale.

\subsection{Funzioni del prodotto}
Le funzionalità principali del sistema si dividono nelle seguenti aree:
\begin{itemize}
    \item \textbf{Gestione Editor visivo}: Area di data entry con supporto al markup Markdown e finestra affiancata per il rendering grafico in tempo reale.
    \item \textbf{Supporto AI e Manipolazione Testo}: Funzioni di riassunto, traduzione\textsubscript{G} multilingua (Inglese, Francese, Spagnolo; Tedesco come estensione facoltativa) e riscrittura del testo per migliorarne forma e scorrevolezza (VE-5.2, VE-1.1).
    \item \textbf{Analisi Critica ("Sei Cappelli")}: Applicazione di prompt\textsubscript{G} specifici per valutare il testo secondo diverse prospettive (imparziale, emozionale, ottimista, critica negativa, creativa, logica).
    \item \textbf{Distant Writing}: Generazione di interi blocchi di testo da parte dell'LLM a partire da un prompt testuale fornito dall'utente.
    \item \textbf{Gestione Dati}: Salvataggio e lettura delle note in file di testo locali (requisito base), con possibile estensione opzionale per il salvataggio strutturato, la modifica retroattiva e l'interconnessione tramite link.
\end{itemize}

\subsection{Caratteristiche degli utenti}
\begin{itemize}
	\item \textbf{Utente Finale}: Un individuo che necessita di supporto nella stesura di testi, nel brainstorming o nello studio. Non è richiesta alcuna competenza tecnica o informatica pregressa: il sistema deve risultare semplice e di utilizzo intuitivo.
\end{itemize}

\subsection{Piattaforma di esecuzione}
Il prodotto è un'applicazione web accessibile tramite browser su sistemi operativi desktop (Windows, macOS, Linux). 

Non è previsto supporto nativo per dispositivi mobile e non sono richiesti requisiti hardware particolari oltre a quelli necessari per l'esecuzione di un browser moderno e per la connessione ai servizi API esterni.

\subsection{Vincoli generali}
I vincoli tecnologici, architetturali, di qualità, di strumento e di rilascio del progetto sono elencati e tracciati in modo puntuale nella Tabella~\ref{tab:req-vincolo} (Requisiti di Vincolo, Sezione~\ref{sec:req-vincolo}), che costituisce l'unica fonte di riferimento per la loro verifica.

In sintesi, il progetto è vincolato da: tecnologie web standard lato client (R1-V-O), con libertà di scelta sul framework\textsubscript{G} front-end adottato (R11-V-O, nel caso del gruppo, Vue.js); interazione con i modelli LLM forniti dall'azienda proponente tramite API conformi allo standard OpenAI (R2-V-O, R3-V-O); adozione di pattern architetturali che favoriscano modularità e revisioni frequenti; pubblicazione di codice e documentazione su repository\textsubscript{G} pubblico (R7-V-O);; consegna dei documenti tecnici ed esterni previsti dal capitolato (R8-V-O, R9-V-O, R10-V-O); distribuzione del prodotto tramite archivio Docker, con deploy su server esterno come estensione opzionale (R12-V-O, R13-V-Op).

Per la modellazione UML\textsubscript{G} e i diagrammi, il team di sviluppo ha scelto di adottare internamente lo strumento \textbf{Lucidchart}.

\subsection{Assunzioni e dipendenze}
Il corretto funzionamento delle funzionalità AI (riassunto, traduzione, riscrittura, Sei Cappelli, Distant Writing) dipende dalla disponibilità dei servizi LLM esterni configurati dall'utente tramite endpoint\textsubscript{G} API compatibile con lo standard OpenAI conseguentemente, il sistema non garantisce il funzionamento di tali funzionalità in assenza di connettività di rete o in caso di indisponibilità del servizio AI configurato.